
\medskip

\emph{Principaux domaines abordés} : Suites numériques.
Algorithmique et programmation.

\bigskip

%On s'intéresse au développement d'une bactérie. 
%
Dans cet exercice, on modélise le développement d'une bactérie avec les hypothèses suivantes : cette bactérie a une probabilité $0,3$ de mourir sans descendance et une probabilité $0,7$ de se diviser en deux bactéries filles.
%
%Dans le cadre de cette expérience, on admet que les lois de reproduction des bactéries sont les mêmes pour toutes les générations de bactéries qu'elles soient mère ou fille.
%
Pour tout entier naturel $n$, on appelle $p_n$ la probabilité d'obtenir au plus $n$ descendances pour une bactérie. 

On admet que, d'après ce modèle, la suite $\left(p_n\right)$ est définie de la façon suivante :

$p_0 = 0,3$ et, pour tout entier naturel $n$, $p_{n+1} = 0,3 + 0,7p_n^2$.

%\begin{minipage}{0.68\linewidth}
\begin{enumerate}
\item La feuille de calcul ci-dessous donne des valeurs approchées de la suite $\left(p_n\right)$
	\begin{enumerate}
		\item %Déterminer les valeurs exactes de $p_1$ et $p_2$ (masquées dans la feuille de calcul) et interpréter ces valeurs dans le contexte de l'énoncé.
$\bullet~~$ $p_1 = 0,3 + 0,7 \times 0,3^2 = 0,3 + 0,7 \times 0,09 = 0,3 + 0,063 = 0,363$ ;

$\bullet~~$ $p_2 = 0,3 + 0,7 \times 0,363^2 = 0,3 + 0,7 \times \np{0,131769} = 0,3 + \np{0,0922383} = \np{0,3922383}$.

		\item %Quelle est la probabilité, arrondie à $10^{-3}$ près, d'obtenir au moins 11 générations de bactéries à partir d'une bactérie de ce type ?
		La probabilité d'obtenir au moins 11 générations de bactéries à partir d'une bactérie de ce type est égale à $1 - p_{10} \approx 1 - \np{0,42802018} \approx \np{0,571980} \approx 0,572$ au millième près.
		
		\item %Formuler des conjectures sur les variations et la convergence de la suite $\left(p_n\right)$.
D'après le tableau la suite $\left(p_n\right)$ semble être croissante et semble aussi avoir une limite puisque les quatre derniers résultats commencent par \np{0,4285}...
	\end{enumerate}

\item 
	\begin{enumerate}
		\item On veut démontrer par récurrence sur $n$ que, pour tout entier naturel $n,\: 0 \leqslant p_n \leqslant p_{n+1} \leqslant0,5$.
		
\emph{Initialisation}  : On a $0 \leqslant 0,3 \leqslant 0,363 \leqslant 0,5$, soit $0 \leqslant p_0 \leqslant p_1 \leqslant 0,5$ : la relation est vraie au rang 0.

\emph{Hérédité} : soit $n \in \N$ et supposons que $0 \leqslant p_n \leqslant p_{n+1} \leqslant 0,5$.

Si ces nombres positifs sont rangés dans dans cet ordre, leurs carrés aussi, soit 

$0 \leqslant p_n^2 \leqslant p_{n+1}^2 \leqslant 0,5^2$, puis en multipliant par 0,7 :

$0 \leqslant 0,7 \times p_n^2 \leqslant 0,7 \times p_{n+1}^2 \leqslant 0,7 \times 0,5^2$ et enfin en ajoutant à chaque membre 0,3 :

$0,3 \leqslant 0,3 + 0,7 \times p_n^2 \leqslant 0,3 + 0,7 \times p_{n+1}^2 \leqslant 0,3 + 0,7 \times 0,5^2$, soit 

$0,3 \leqslant p_{n+1} \leqslant p_{n+2} \leqslant 0,475$. On peut donc écrire :

$0 \leqslant p_{n+1} \leqslant p_{n+2} \leqslant 0,5$ : la relation est vraie au rang $n + 1$.

\emph{Conclusion} : la relation est vraie au rang 0 et si elle vraie au rang $n \in \N$, elle est encore vraie au rang $n + 1$ : d'après la principe de récurrence, pour tout entier naturel $n$, on a  $0 \leqslant p_n \leqslant p_{n+1} \leqslant0,5$
		\item %Justifier que la suite $\left(p_n\right)$ est convergente.
Le résultat précédent montre que la suite $\left(p_n\right)$ est croissante et majorée par 0,5 : elle est donc convergente vers une limite $L$ telle que $L \leqslant 0,5$.
	\end{enumerate}
\item %On appelle $L$ la limite de la suite $\left(p_n\right)$.
	\begin{enumerate}
		\item %Justifier que $L$ est solution de l'équation 
		
%		\[0,7x^2 - x  + 0,3 = 0\]
Puisque $\displaystyle\lim_{n \to + \infty} p_n = \displaystyle\lim_{n \to + \infty} p_{n+1} = L$, la relation de récurrence donne :

\[L = 0,3 + 0,7L^2 \iff 0,7L^2 - L + 0,3 = 0\]

$L$ est solution de l'équation $0,7x^2 - x  + 0,3 = 0$.
		\item %Déterminer alors la limite de la suite $\left(p_n\right)$.
		On a $\Delta = 1 - 4 \times 0,7 \times 0,3 = 1 - 0,84 = 0,16 = 0,4^2 > 0$. Il y a donc deux solutions :
		
\[L_1 = \dfrac{1 + 0,4}{2\times 0,7} = 1, \quad \text{et}\quad L_2 = \dfrac{1 - 0,4}{2\times 0,7} = \frac{0,6}{1,4} = \dfrac{6}{14} = \dfrac37 \approx \np{0,43}\]
		
On ne peut retenir la solution $L_1$ puisque $\left(p_n\right)$ est majorée par 0,5. Il reste donc $L_2$ : 
$\displaystyle\lim_{n \to + \infty} p_n = \dfrac37$.
	\end{enumerate}
\end{enumerate}

\begin{enumerate}[resume]
\item La fonction suivante, écrite en langage Python, a pour objectif de renvoyer les $n$ premiers termes de la suite $\left(p_n\right)$.

On complète, les lignes 2, 4 et 5 de cette fonction de façon à ce que la fonction \texttt{suite (n)} retourne, sous forme de liste, les $n$ premiers termes de la suite.

\begin{center}
\begin{tabularx}{0.4\linewidth}{c |l|}\cline{2-2}
1 &\texttt{\textbf{def} suite(n) :}\\
2 &\quad \texttt{p = {\blue 0.3}}\\
3 &\quad \texttt{s = [p]}\\
4 &\quad \texttt{for i in range ({\blue n - 1})}\\
5 &\quad \qquad \texttt{p = {\blue 0.3+0.7*p**2}}\\
6 &\quad \qquad \texttt{s.append(p)}\\
7 &\quad \texttt{\textbf{return} (s)}\\ \cline{2-2}
\end{tabularx}
\end{center}

%Recopier, sur votre copie, cette fonction en complétant les lignes 2, 4 et 5 de façon à ce que la fonction \texttt{suite (n)} retourne, sous forme de liste, les $n$ premiers termes de la suite.
\end{enumerate}
